\section{Full per-family results}
\label{app:full}

\Cref{tab:appendix} reports per-family, per-seed means for the two trained arms.
% Table 4: appendix per-family comparison (generated by gen_tables.py; seed-avg of per-task metrics)
\begin{table}[t]
\centering
\caption{Per-family, per-seed means for the two trained arms (task-macro of per-task seed means). \#1 uses per-family early-stop checkpoints (37-task slice; families as in the training scheduler), \#2 the full-run checkpoint. DeepLoc (multi\_noul) and double-sequence tasks are outside this harness.}
\label{tab:appendix}
\footnotesize
\begin{tabular}{llccc}
\toprule
Family & Primitive & \#1 (mean$\pm$std over 3 seeds) & \#2 & n \\ \midrule
DeepSTARR & score & 0.035$\pm$0.027 & 0.053$\pm$0.023 & 3 \\
GUE & choice & 0.464$\pm$0.017 & 0.478$\pm$0.013 & 3 \\
GenomicBenchmarks & choice & 0.402$\pm$0.082 & 0.305$\pm$0.014 & 3 \\
GenomicBenchmarks & noul & 0.478$\pm$0.017 & 0.629$\pm$0.005 & 3 \\
ProteinGym & score & 0.015$\pm$0.058 & 0.022$\pm$0.015 & 3 \\
TAPE & score & -0.014$\pm$0.047 & -0.027$\pm$0.043 & 3 \\
dnagpt\_pools & choice & 0.519$\pm$0.006 & 0.514$\pm$0.030 & 3 \\
local\_snapshots & choice & 0.534$\pm$0.016 & 0.536$\pm$0.030 & 3 \\
\bottomrule
\end{tabular}
\end{table}


\section{Adaptation cost and ablation tables}

% Table 3: adaptation cost (generated by gen_tables.py)
\begin{table}[t]
\centering
\caption{Adaptation cost on 11 held-out tasks: \#1 zero-shot (no labels, no training) vs \#2 new linear head fitted on $N$ labels (frozen shared encoder). \#2 catches or beats \#1's zero-shot on 7/11 tasks already at $N{=}32$.}
\label{tab:adapt}
\footnotesize
\begin{tabular}{llcccc}
\toprule
Task & Metric & \#1 zero-shot & \#2 $N{=}32$ & $N{=}128$ & $N{=}512$ \\
\midrule
gue\_H3 & accuracy & 0.555 & 0.510 & 0.510 & -- \\
gue\_H3K14ac & accuracy & 0.445 & 0.543 & 0.543 & -- \\
gue\_H3K4me1 & accuracy & 0.430 & 0.532 & 0.532 & -- \\
gue\_H3K4me2 & accuracy & 0.450 & 0.541 & 0.541 & -- \\
gue\_H3K79me3 & accuracy & 0.490 & 0.492 & 0.492 & -- \\
gue\_H3K9ac & accuracy & 0.485 & 0.451 & 0.451 & -- \\
pg\_ACE2\_HUMAN\_Chan\_2020 & spearman & -0.176 & 0.029 & -0.005 & 0.040 \\
pg\_B2L11\_HUMAN\_Dutta\_2010\_binding-Mcl-1 & spearman & 0.133 & -0.054 & -- & -- \\
pg\_CATR\_CHLRE\_Tsuboyama\_2023\_2AMI & spearman & 0.066 & -0.086 & -0.039 & 0.135 \\
pg\_CSN4\_MOUSE\_Tsuboyama\_2023\_1UFM & spearman & -0.034 & 0.171 & 0.101 & -0.088 \\
pg\_F7YBW8\_MESOW\_Ding\_2023 & spearman & -0.112 & 0.416 & -- & -- \\
\bottomrule
\end{tabular}
\end{table}



\begin{figure*}[t]
\centering
\includegraphics[width=0.98\textwidth]{fig4_ablation.pdf}
\caption{\textbf{Controlled mechanism ablations.}
(a) Noul candidate phrasing on three seed-averaged stable checkpoints (same
weights; only candidate text varies). No phrasing reaches the per-task-head
reference (0.629, dashed); the best (``false: no / true: yes'', $0.553\pm0.023$)
beats the \emph{trained} phrasing (hatched) on all three seeds --- reversing the
single-checkpoint finding, i.e., inference-level ablations are also
checkpoint-dependent.
(b) Frozen VLM on synthetic property maps (multimodal route-A pilot): plain
candidate likelihood collapses to one class (acc 0.040, 1/7 classes used);
contrastive decoding removes the collapse (0.140, 4/7) but remains at chance ---
decoding fixed, signal absent.}
\label{fig:ablation}
\end{figure*}

\label{app:ablations}

% Table 5: noul phrasing x 3 stable checkpoints (gen_tables.py)
\begin{table}[t]
\centering
\caption{Noul AUROC on GB tasks for six candidate phrasings, evaluated on three seed-averaged stable checkpoints (same weights per column; only rendered candidate text varies). No phrasing reaches the per-task-head reference (0.629). The trained phrasing is marked \dag.}
\label{tab:phrasing}
\small
\begin{tabular}{lcccc}
\toprule
Phrasing & s1 & s2 & s3 & mean$\pm$std \\ \midrule
false: no / true: yes & 0.572 & 0.568 & 0.521 & 0.553$\pm$0.023 \\
false / true & 0.566 & 0.523 & 0.541 & 0.543$\pm$0.018 \\
proposition (``yes, it does'') & 0.517 & 0.513 & 0.549 & 0.526$\pm$0.016 \\
no / yes (trained)$^\dagger$ & 0.477 & 0.536 & 0.510 & 0.508$\pm$0.024 \\
descriptive & 0.527 & 0.506 & 0.431 & 0.488$\pm$0.041 \\
neg / pos words & 0.513 & 0.468 & 0.472 & 0.484$\pm$0.020 \\
\bottomrule
\end{tabular}
\end{table}

% Table 6: TAPE cross-family score arm, 3 seeds (gen_tables.py)
\begin{table}[t]
\centering
\caption{Cross-family score transfer (TAPE held out; DeepSTARR and ProteinGym score families in training). Spearman per seed; ``abs'' = trained absolute value. Stability's positive $\Delta$ reflects a pathological frozen baseline ($-0.218$) regressing toward zero, not genuine transfer.}
\label{tab:tape}
\small
\begin{tabular}{llccccc}
\toprule
Task & & frozen & tr.\ s1 & tr.\ s2 & tr.\ s3 & $\Delta$ mean$\pm$std \\ \midrule
tape\_fluorescence & abs & 0.037 & -0.008 & 0.026 & -0.075 & -0.056$\pm$0.042 \\
tape\_stability & abs & -0.218 & -0.018 & 0.020 & -0.179 & +0.159$\pm$0.086 \\
\bottomrule
\end{tabular}
\end{table}


\section{Multimodal route-A pilot}
\label{app:vlm}

Protein sequences without resolved structures were rendered as synthetic
amino-acid property heatmaps (7 normalized physicochemical tracks $\times$ position,
matplotlib viridis) and presented to a frozen Qwen2.5-VL-3B \citep{bai2025qwen25vl}
alongside the fold-class question. Plain full-candidate likelihood collapsed to a single
candidate on all 200 images (accuracy 0.040; macro-F1 0.011; 1/7 classes predicted) ---
candidate text priors dominate and the image does not change the ranking. Contrastive
decoding, $\log p(\text{cand}\mid\text{img},q)-\log p(\text{cand}\mid q)$, restores
multi-class behavior (4/7 classes, accuracy 0.140) but sits at chance
($1/7=0.143$), far below the majority baseline (0.325) and the text-only frozen LM
(0.320). Image sanity checks rule out rendering failure (per-image statistics and
cross-image correlations normal). Conclusion: route A (synthetic renders) with a frozen
VLM is a dead end; the reusable finding is that \emph{frozen-VLM candidate-likelihood
evaluation should default to contrastive decoding}. Real-image localization (HPA) and
fine-tuned VLM arms are a separate work.

\section{Data, licensing, and disclosures}

\paragraph{Data and licensing.}
All sources are public datasets with license, revision and SHA-256 recorded in download
receipts; no human-subject data. RNAcompete is the 14-RBP test subset (not the full
${\sim}$200); DeepGOZero/PEER/DeepSEA remain in the catalog as needs-upstream-work and are
not claimed as coverage. The splice label-convention doubt (\cref{sec:admission}) is
explicitly recorded for independent verification.

\paragraph{LLM-use disclosure.}
[To be finalized per venue checklist: writing assistance disclosure.]
\section{Reproducibility}
\label{app:repro}

Every figure and table regenerates from on-disk JSONs via one script per artifact
(\texttt{paper\_v2/figures/gen\_*.py}); no metric is hardcoded in the plotting or table
code. Each run persists \texttt{run\_config.json} (flags, seeds, caps),
\texttt{train\_trace.jsonl} (per-step loss/grad) and \texttt{dev\_curve.jsonl} (per-family
dev scores); evaluation subsets are drawn with a fixed sampling seed (20261001) rather than
first-$N$ --- a bug we found and fixed: class-ordered test files plus first-$N$ sampling
produced a spurious noul AUROC of 0.387 that looked like a polarity bug but was a sampling
artifact. Hardware: single H20/A800-class GPU per run; nondeterminism is absorbed by the
seed protocol of \cref{sec:seedavg}.
